% =============================================================================
%  Acelerador em FPGA do Supervisor de Seguranca LiDAR de uma Cadeira de Rodas
%  Ponto de Controle 2 (PC2) -- Implementacao RTL e Verificacao por Equivalencia
%  PCR / Co-Projeto Hardware-Software -- UnB 2026/1
%
%  Compilacao: pdflatex relatorio_pcr_pc2.tex  (2x, por causa das referencias)
%  Requer os .dat em figs/ (gerados por golden/gen_figs.py e gen_figs_cmp.py).
%
%  Este e o SEGUNDO ponto de controle: o RTL esta escrito e verificado por
%  simulacao (equivalencia bit a bit contra o modelo de ouro, sobre varreduras
%  reais). A lat encia vem da simulacao GHDL. A SINTESE (recursos, timing,
%  potencia, block design) ainda NAO foi rodada -- e o objeto de PC3.
% =============================================================================
\documentclass[conference]{IEEEtran}
\IEEEoverridecommandlockouts

\usepackage[utf8]{inputenc}
\usepackage[T1]{fontenc}
\usepackage[brazilian]{babel}
\usepackage{graphicx}
\usepackage{amsmath,amssymb}
\usepackage{booktabs}
\usepackage{cite}
\usepackage{url}
\usepackage{xcolor}
\usepackage{tikz}
\usepackage{pgfplots}
\usepackage{listings}
\usepackage{siunitx}
\pgfplotsset{compat=1.17}
\usetikzlibrary{arrows.meta, positioning, fit, calc}

\sisetup{output-decimal-marker={,}, per-mode=symbol, detect-all}

\lstset{basicstyle=\ttfamily\scriptsize, breaklines=true, columns=fullflexible,
        frame=single, framesep=3pt, xleftmargin=2pt}

\begin{document}

\title{Acelerador em FPGA para o Supervisor de Seguran\c{c}a\\
LiDAR de uma Cadeira de Rodas Semi-Assistida:\\
\large Ponto de Controle 2 --- Implementa\c{c}\~ao RTL e Verifica\c{c}\~ao por Equival\^encia}

\author{
\IEEEauthorblockN{Vítor Gonçalves de Andrade Silva}
\IEEEauthorblockA{Engenharia Eletrônica --- FCTE/UnB\\
Email: 261025182@aluno.unb.br}
\and
\IEEEauthorblockN{Thiago Henrique Oliveira Souza}
\IEEEauthorblockA{Engenharia Eletrônica --- FCTE/UnB\\
Email: 180131672@aluno.unb.br}
}

\maketitle

% =============================================================================
\begin{abstract}
Este segundo ponto de controle implementa em VHDL a arquitetura proposta em PC1
--- um CORDIC pipelinado de \num{16} est\'agios alimentando \num{19} pistas
paralelas de coeficiente constante --- e a verifica por \textbf{equival\^encia bit
a bit} contra um modelo de refer\^encia, sobre tr\^es varreduras \emph{reais}
capturadas na cadeira. As tr\^es passam com zero diverg\^encia. A
verifica\c{c}\~ao exp\^os o risco antecipado em PC1: um \textbf{defeito de faixa}
no formato Q3.13. O \^angulo do feixe atinge \SI{-5.85}{\radian} e n\~ao cabe em
$\pm\SI{4.0}{\radian}$, saturando e projetando o ponto no lugar errado (erro de
\SI{402}{\milli\metre}); envolvendo o \^angulo em $[-\pi,\pi)$ antes da
quantiza\c{c}\~ao, o erro cai a \SI{0.99}{\milli\metre} m\'edio --- cerca de um
LSB. A lat\^encia, medida na simula\c{c}\~ao GHDL, \'e de \num{751} ciclos
(\SI{7.51}{\micro\second} a \SI{100}{\mega\hertz}); contra o \emph{baseline}
medido na pr\'opria Raspberry~Pi~5 (\SI{3.54}{\milli\second} escalar,
\SI{0.488}{\milli\second} vetorizado), \num{472}$\times$ e \num{65}$\times$. A
s\'intese --- recursos, \emph{timing}, pot\^encia e integra\c{c}\~ao em SoC ---
ainda \textbf{n\~ao foi executada} e \'e o objeto do ponto de controle final.
\end{abstract}

% =============================================================================
\section{Recapitula\c{c}\~ao}

Em PC1 identificamos, por perfilamento, o teste de corredor do supervisor de
seguran\c{c}a como o bloco a acelerar: \'e o n\'o mais pesado
(\SI{17.7}{\percent} da CPU da Pi~5) e o \'unico cuja lat\^encia tem
consequ\^encia f\'isica, pois s\'o corta o avan\c{c}o e a cadeira ainda percorre
dist\^ancia por in\'ercia~\cite{fehr2000}. O algoritmo projeta os \num{720}
feixes no referencial \texttt{base\_link} e testa, para \num{19} dire\c{c}\~oes
candidatas $\delta_k$, se cada ponto cai no corredor f\'isico da cadeira
(largura $W=\SI{0.61}{\metre}$, frente $X_f=\SI{0.667}{\metre}$) --- \num{9918}
avalia\c{c}\~oes por varredura, custo $O(N\cdot K)$. Propusemos migrar esse custo
para hardware e manter a fun\c{c}\~ao custo $O(K)$ no ARM, onde vivem os pesos
ajustados em campo. Este documento implementa e verifica; a s\'intese fica para
PC3.

% =============================================================================
\section{Implementa\c{c}\~ao RTL}

O RTL est\'a organizado em quatro m\'odulos VHDL (Tab.~\ref{tab:rtl}), do CORDIC
ao envelope AXI, todos parametriz\'aveis pelos formatos fixados em PC1.

\begin{table}[!t]
\caption{M\'odulos RTL implementados}
\label{tab:rtl}
\centering\footnotesize
\begin{tabular}{@{}ll@{}}
\toprule
\textbf{Arquivo} & \textbf{Fun\c{c}\~ao} \\
\midrule
\texttt{cordic\_pkg.vhd}   & constantes (\textbf{geradas} pelo modelo de ouro) \\
\texttt{cordic.vhd}        & CORDIC rota\c{c}\~ao, \num{16} est\'agios, \num{0} DSP \\
\texttt{lane.vhd}          & uma pista: rota\c{c}\~ao de coef.\ constante $+$ teste \\
\texttt{corridor\_core.vhd}& proje\c{c}\~ao $+$ \num{19} pistas $+$ redu\c{c}\~ao de m\'inimo \\
\texttt{lidar\_accel\_axi.vhd} & envelope AXI4-Stream (feixes) $+$ AXI4-Lite (ctrl/folgas) \\
\bottomrule
\end{tabular}
\end{table}

\paragraph{CORDIC}
Implementado em modo rota\c{c}\~ao~\cite{volder1959}, pipelinado em \num{16}
est\'agios, usando apenas somadores e deslocamentos cabeados --- sem nenhum
multiplicador. Como o la\c{c}o s\'o converge para $|z|\le\SI{1.7434}{\radian}$, um
est\'agio de redu\c{c}\~ao de quadrante traz $\theta$ para $[-\pi/2,\pi/2]$ e
inverte o sinal na sa\'ida. A op\c{c}\~ao por CORDIC, e n\~ao por LUT de seno, se
justifica porque \texttt{angle\_min}, o incremento e o \emph{yaw} s\~ao
par\^ametros de calibra\c{c}\~ao reescritos em tempo de execu\c{c}\~ao via
AXI4-Lite --- uma ROM indexada pelo \'indice do feixe exigiria fix\'a-los em
s\'intese.

\paragraph{Pistas}
Em cada pista, $\cos\delta_k$ e $\sin\delta_k$ s\~ao \emph{generics} constantes de
s\'intese; a rota\c{c}\~ao vira quatro multiplica\c{c}\~oes de coeficiente fixo. O
\emph{core} instancia as \num{19} pistas em paralelo, de modo que a vaz\~ao alvo
de \textbf{1 feixe/ciclo} se mant\'em: cada feixe projetado \'e avaliado
simultaneamente pelos \num{19} candidatos.

\paragraph{Constantes geradas, n\~ao transcritas}
As tabelas de \texttt{atan}, o ganho do CORDIC e os coeficientes das pistas em
\texttt{cordic\_pkg.vhd} s\~ao \textbf{gerados} por \texttt{golden/gen\_pkg.py} a
partir do mesmo modelo que serve de refer\^encia. Assim, RTL e refer\^encia casam
\emph{por constru\c{c}\~ao}: um erro de digita\c{c}\~ao numa tabela de constantes
seria, de outro modo, indetect\'avel por inspe\c{c}\~ao.

% =============================================================================
\section{Modelo de Ouro e Metodologia de Verifica\c{c}\~ao}

Conforme planejado em PC1, a verifica\c{c}\~ao \'e por \textbf{equival\^encia bit
a bit}, e n\~ao por inspe\c{c}\~ao de formas de onda:
\begin{enumerate}
\item um \textbf{modelo em ponto fixo} (\texttt{golden/model.py}) emula o RTL
exatamente --- mesmo CORDIC, mesmos deslocamentos, mesmos truncamentos;
\item um \textbf{modelo em ponto flutuante} serve de refer\^encia f\'isica, para
medir o erro de quantiza\c{c}\~ao;
\item o modelo gera o \textbf{est\'imulo} e o \textbf{esperado}; o testbench VHDL
(\texttt{tb/tb\_corridor\_core.vhd}) l\^e os arquivos, injeta os \num{720} feixes
e compara as \num{19} folgas;
\item o est\'imulo n\~ao \'e sint\'etico: s\~ao tr\^es varreduras \textbf{reais}
da cadeira, capturadas com um balde posicionado \SI{21}{\centi\metre} \`a direita
do eixo --- exatamente o antigo ponto cego (Fig.~\ref{fig:corredor}).
\end{enumerate}

\begin{figure}[!t]
\centering
\begin{tikzpicture}
\begin{axis}[width=\columnwidth, height=5.2cm,
  xlabel={$X$ (m) --- \`a frente do eixo traseiro},
  ylabel={$Y$ (m)}, xmin=0, xmax=2.6, ymin=-1.1, ymax=1.5,
  legend pos=north east, legend style={font=\tiny},
  grid=major, grid style={gray!20}, tick label style={font=\scriptsize},
  label style={font=\scriptsize}, axis equal image=false]
  \addplot[only marks, mark=*, mark size=0.5pt, gray!60]
    table[x=x, y=y] {figs/scan_out.dat};
  \addlegendentry{fora do corredor}
  \addplot[only marks, mark=*, mark size=0.9pt, red!80!black]
    table[x=x, y=y] {figs/scan_in.dat};
  \addlegendentry{obst\'aculo (no corredor)}
  \draw[blue!60!black, thick] (axis cs:0.687,-0.305) rectangle (axis cs:2.6,0.305);
  \node[blue!60!black, font=\tiny] at (axis cs:2.1,0.42) {corredor da cadeira};
  \draw[fill=black] (axis cs:0.667,0.336) circle (1.6pt);
  \node[font=\tiny, anchor=south] at (axis cs:0.667,0.38) {LiDAR};
\end{axis}
\end{tikzpicture}
\caption{Varredura real projetada no \texttt{base\_link} (est\'imulo do
testbench). O LiDAR est\'a \SI{0.336}{\metre} \`a esquerda do eixo. Os pontos em
vermelho caem no corredor que a cadeira ocupar\'a --- entre eles o balde, que o
cone antigo n\~ao enxergava.}
\label{fig:corredor}
\end{figure}

% =============================================================================
\section{Resultados de Simula\c{c}\~ao}

\subsection{Equival\^encia}
As tr\^es varreduras reais passam pelo testbench com \textbf{zero diverg\^encia}:
para cada uma, as \num{19} folgas produzidas pelo RTL simulado coincidem com o
esperado gerado pelo modelo de ponto fixo. A Fig.~\ref{fig:swhw} confronta, por
candidato, a folga do software (ponto flutuante, a refer\^encia f\'isica) com a
sa\'ida efetiva da simula\c{c}\~ao VHDL; as curvas se sobrep\~oem.

\begin{figure}[!t]
\centering
\begin{tikzpicture}
\begin{axis}[width=\columnwidth, height=4.4cm,
  xlabel={candidato de desvio $\delta_k$ (graus)},
  ylabel={folga (m)}, xmin=-48, xmax=48, xtick={-45,-30,-15,0,15,30,45},
  legend pos=north west, legend style={font=\scriptsize},
  grid=major, grid style={gray!20}, tick label style={font=\scriptsize},
  label style={font=\scriptsize}]
  \addplot[thick, blue!70!black, mark=o, mark size=2pt]
    table[x=delta, y=sw] {figs/sw_hw_scan0.dat};
  \addlegendentry{SW (ponto flutuante)}
  \addplot[thick, red, only marks, mark=x, mark size=3pt]
    table[x=delta, y=hw] {figs/sw_hw_scan0.dat};
  \addlegendentry{HW (RTL simulado)}
\end{axis}
\end{tikzpicture}
\caption{Folga por candidato: software \emph{versus} sa\'ida real do RTL, para a
varredura de teste. As curvas se sobrep\~oem em toda a faixa --- o desvio \'e
invis\'ivel nesta escala, precisamente o resultado desejado.}
\label{fig:swhw}
\end{figure}

\subsection{Precis\~ao: um defeito de formato}
\label{sec:prec}
A an\'alise de precis\~ao confirmou o risco levantado em PC1 e produziu o achado
mais instrutivo desta fase. O primeiro modelo em ponto fixo divergia da
refer\^encia em at\'e \SI{402}{\milli\metre} --- inaceit\'avel num sistema de
seguran\c{c}a, e grande demais para ser arredondamento.

A causa n\~ao era ru\'ido de quantiza\c{c}\~ao, e sim \textbf{estouro de faixa}. O
\^angulo do feixe no referencial do sensor,
$\theta_i = a_{\min} + i\cdot\Delta a + \theta_L$, atinge \SI{-5.85}{\radian}; o
formato Q3.13 comporta apenas $\pm\SI{4.0}{\radian}$. O valor \textbf{saturava}, e
o feixe era projetado numa dire\c{c}\~ao completamente errada. Envolvendo
$\theta$ em $[-\pi,\pi)$ antes da quantiza\c{c}\~ao --- opera\c{c}\~ao barata,
feita no host --- o erro desaba (Tab.~\ref{tab:prec}).

\begin{table}[!t]
\caption{Erro do ponto fixo contra a refer\^encia em ponto flutuante}
\label{tab:prec}
\centering\footnotesize
\begin{tabular}{@{}lrr@{}}
\toprule
 & \textbf{erro m\'edio} & \textbf{erro m\'aximo} \\
\midrule
$\theta$ saturado (defeito) & \SI{17.89}{\milli\metre} & \SI{402.15}{\milli\metre} \\
$\theta$ envolvido (corrigido) & \textbf{\SI{0.99}{\milli\metre}} & \textbf{\SI{5.27}{\milli\metre}} \\
\bottomrule
\end{tabular}
\end{table}

O erro final \'e da ordem de \textbf{um LSB} de Q6.10 (\SI{0.98}{\milli\metre}) e
representa \SI{0.9}{\percent} da zona de parada de \SI{600}{\milli\metre}: a
quantiza\c{c}\~ao \emph{n\~ao} \'e o fator limitante. A li\c{c}\~ao \'e que, em
ponto fixo, a escolha da \textbf{faixa} costuma ser mais cr\'itica que a da
\textbf{resolu\c{c}\~ao} --- perdemos \num{402} milímetros por tr\^es bits de
parte inteira, n\~ao por falta de bits fracion\'arios. A Fig.~\ref{fig:erro}
mostra o res\'iduo com sinal nas tr\^es varreduras.

\begin{figure}[!t]
\centering
\begin{tikzpicture}
\begin{axis}[width=\columnwidth, height=4.2cm,
  xlabel={candidato de desvio $\delta_k$ (graus)},
  ylabel={erro HW$-$SW (mm)}, xmin=-48, xmax=48,
  xtick={-45,-30,-15,0,15,30,45},
  legend pos=south west, legend style={font=\tiny},
  grid=major, grid style={gray!20}, tick label style={font=\scriptsize},
  label style={font=\scriptsize}]
  \draw[gray, dashed] (axis cs:-48,0.98) -- (axis cs:48,0.98);
  \draw[gray, dashed] (axis cs:-48,-0.98) -- (axis cs:48,-0.98);
  \node[gray, font=\tiny, anchor=west] at (axis cs:-46,1.7) {$\pm 1$ LSB};
  \addplot[only marks, mark=*, mark size=1.6pt, blue!70!black]
    table[x=delta, y=err_mm] {figs/err_scan0.dat};
  \addlegendentry{varredura 0}
  \addplot[only marks, mark=square*, mark size=1.6pt, red!70!black]
    table[x=delta, y=err_mm] {figs/err_scan1.dat};
  \addlegendentry{varredura 1}
  \addplot[only marks, mark=triangle*, mark size=2pt, orange!90!black]
    table[x=delta, y=err_mm] {figs/err_scan2.dat};
  \addlegendentry{varredura 2}
\end{axis}
\end{tikzpicture}
\caption{Res\'iduo HW$-$SW, com sinal, nas tr\^es varreduras (\num{57}
compara\c{c}\~oes). Linhas tracejadas: $\pm1$ LSB de Q6.10
(\SI{0.98}{\milli\metre}). O erro m\'aximo, \SI{5.27}{\milli\metre}, ocorre quando
a quantiza\c{c}\~ao faz um ponto pr\'oximo da borda cruzar o limiar
$|y|\le W/2$ --- efeito de um crit\'erio por limiar, n\~ao ac\'umulo aritm\'etico.}
\label{fig:erro}
\end{figure}

\subsection{Lat\^encia medida em simula\c{c}\~ao}
A lat\^encia \'e o n\'umero de ciclos \textbf{contado na simula\c{c}\~ao} GHDL,
n\~ao uma estimativa anal\'itica: \num{751} ciclos por varredura, ou
\SI{7.51}{\micro\second} a \SI{100}{\mega\hertz} --- pr\'oximo da previs\~ao de
PC1 ($\sim N$ mais a profundidade do \emph{pipeline}). O \emph{baseline} de
software foi medido na pr\'opria Raspberry~Pi~5 que embarca o sistema
(Tab.~\ref{tab:speedup}).

\begin{table}[!t]
\caption{Tempo por varredura e ganho (lat\^encia HW por simula\c{c}\~ao)}
\label{tab:speedup}
\centering\footnotesize
\begin{tabular}{@{}lrr@{}}
\toprule
\textbf{Implementa\c{c}\~ao} & \textbf{Tempo} & \textbf{Speedup} \\
\midrule
Software escalar (Pi 5)      & \SI{3.54}{\milli\second}   & --- \\
Software vetorizado (NumPy)  & \SI{0.488}{\milli\second}  & \num{7.3}$\times$ \\
\textbf{HW \SI{100}{\mega\hertz}} (\num{751} ciclos) & \textbf{\SI{7.51}{\micro\second}} & \num{472}$\times$ / \num{65}$\times$ \\
\bottomrule
\end{tabular}
\end{table}

Vale a interpreta\c{c}\~ao honesta, j\'a poss\'ivel nesta fase: o LiDAR entrega
uma varredura a cada \SI{100}{\milli\second}, e mesmo o software vetorizado
consome apenas \SI{0.5}{\percent} desse or\c{c}amento. O acelerador n\~ao
``salva'' o sistema de perder prazos; seu valor \'e liberar CPU numa Pi que roda
SLAM e EKF e, sobretudo, tornar a lat\^encia \textbf{determin\'istica}. A
quantifica\c{c}\~ao definitiva desse ganho depende da s\'intese (pot\^encia,
$F_{\max}$), ainda pendente.

% =============================================================================
\section{Estado da S\'intese --- Pendente}

A s\'intese ainda \textbf{n\~ao foi executada}. Mantemos, at\'e PC3, apenas a
estimativa anal\'itica de PC1: \num{78} DSP48 ($19\times4$ das pistas mais
\num{2} da proje\c{c}\~ao) e \num{0} BRAM, por ser \emph{streaming}. Esses
n\'umeros de recursos, o fechamento de \emph{timing} (WNS, $F_{\max}$), a
pot\^encia e a integra\c{c}\~ao em \emph{block design} com o ARM e o AXI-DMA
constituem o escopo do ponto de controle final. Registramos a previs\~ao aqui
justamente para confront\'a-la, em PC3, com o que a ferramenta efetivamente
produzir.

\begin{table}[!t]
\caption{O que fecha em PC3}
\label{tab:pc3}
\centering\footnotesize
\begin{tabular}{@{}ll@{}}
\toprule
\textbf{Item} & \textbf{Estado ap\'os PC2} \\
\midrule
RTL (CORDIC, pista, \emph{core}, AXI) & \textbf{implementado} \\
Verifica\c{c}\~ao por equival\^encia    & \textbf{passa (3 varreduras reais)} \\
Defeito de faixa Q3.13                 & \textbf{encontrado e corrigido} \\
Lat\^encia (simula\c{c}\~ao)           & \textbf{\num{751} ciclos} \\
Recursos, \emph{timing}, pot\^encia    & pendente (s\'intese em PC3) \\
\emph{Block design} / grava\c{c}\~ao em placa & pendente (PC3) \\
\bottomrule
\end{tabular}
\end{table}

% =============================================================================
\section{Pr\'oximos Passos (PC3)}
\begin{enumerate}
\item \textbf{Sintetizar} no \texttt{xc7z020clg484-1} (ZedBoard) e extrair
recursos, WNS/$F_{\max}$ e pot\^encia; confrontar os DSP obtidos com a previs\~ao
de \num{78}.
\item \textbf{Integrar em SoC}: empacotar o acelerador como IP e mont\'a-lo com o
ARM Cortex-A9 e o AXI-DMA num \emph{block design}.
\item \textbf{Gravar em placa} e exercitar o driver do ARM
(\texttt{sw/lidar\_accel.c}) sobre o AXI-DMA, fechando o co-projeto.
\end{enumerate}

% =============================================================================
\section{Conclus\~ao Parcial}

O RTL da arquitetura proposta em PC1 est\'a escrito e \textbf{verificado por
equival\^encia bit a bit} contra um modelo de refer\^encia, sobre dados reais da
cadeira: as tr\^es varreduras passam. A verifica\c{c}\~ao confirmou o risco que
hav\'iamos sinalizado --- a faixa do formato de \^angulo --- e o converteu no
resultado mais transfer\'ivel desta fase: um \^angulo de \SI{-5.85}{\radian}
saturando em Q3.13 produziu \SI{402}{\milli\metre} de erro \emph{silencioso}, que
s\'o a compara\c{c}\~ao com uma refer\^encia f\'isica revelou. A lat\^encia
medida em simula\c{c}\~ao (\num{751} ciclos) confirma a ordem de grandeza
prevista. Resta a s\'intese --- recursos, \emph{timing}, pot\^encia e
integra\c{c}\~ao em SoC ---, objeto do ponto de controle final.

% =============================================================================
\begin{thebibliography}{9}

\bibitem{fehr2000}
L.~Fehr, W.~E. Langbein e S.~B. Skaar, ``Adequacy of power wheelchair control
interfaces for persons with severe disabilities: A clinical survey,''
\emph{J. Rehabil. Res. Dev.}, vol.~37, no.~3, pp.~353--360, 2000.

\bibitem{borenstein1991}
J.~Borenstein e Y.~Koren, ``The Vector Field Histogram --- Fast Obstacle
Avoidance for Mobile Robots,'' \emph{IEEE Trans. Robot. Autom.}, vol.~7, no.~3,
pp.~278--288, 1991.

\bibitem{volder1959}
J.~E. Volder, ``The CORDIC Trigonometric Computing Technique,'' \emph{IRE Trans.
Electron. Comput.}, vol.~EC-8, no.~3, pp.~330--334, 1959.

\bibitem{andraka1998}
R.~Andraka, ``A survey of CORDIC algorithms for FPGA based computers,'' em
\emph{Proc. ACM/SIGDA Int. Symp. FPGAs}, 1998, pp.~191--200.

\bibitem{zynqbook}
L.~H. Crockett, R.~A. Elliot, M.~A. Enderwitz e R.~W. Stewart, \emph{The Zynq
Book: Embedded Processing with the ARM Cortex-A9 on the Xilinx Zynq-7000}.
Glasgow: Strathclyde Academic Media, 2014.

\bibitem{ros2docs}
Open Robotics, ``ROS\,2 Jazzy Jalisco Documentation,'' 2024. [Online].
Dispon\'ivel: \url{https://docs.ros.org/en/jazzy/}

\end{thebibliography}

\end{document}
